\section{Fundamentação teórica}\label{sec:fundamentacao}

\subsection{Gestão de resíduos sólidos urbanos no Brasil}

A Lei nº 12.305/2010 instituiu a Política Nacional de Resíduos Sólidos (PNRS),
que define responsabilidades dos titulares do serviço público de limpeza
urbana, exige planos de gestão e estabelece a ordem de prioridade entre não
geração, redução, reutilização, reciclagem, tratamento e disposição final
ambientalmente adequada \cite{brasil2010lei}. A Lei nº 14.026/2020 atualizou o
marco legal do saneamento básico, que inclui o manejo de resíduos sólidos
\cite{brasil2020lei}. Os dados nacionais do setor passaram a ser publicados
pelo Sistema Nacional de Informações em Saneamento Básico (SINISA) e pelo
SINIR \cite{brasil2025sinisa,sinirsdportal}.

Este trabalho trata da etapa de coleta, e não da destinação. A coleta
domiciliar é fortemente influenciada pela geração per capita, pela densidade
aparente dos resíduos e pelo tipo de acondicionamento. O manual do IBAM indica
\num{230} kg/m³ como peso específico aparente do lixo domiciliar solto, na
falta de dados locais, e descreve contêineres plásticos de 120 a \num{1100}~L
e caminhões compactadores de 6 a 19~m³ \cite{monteiro2001manual}. A agência
ambiental dos Estados Unidos reporta faixas de 148 a 178~kg/m³ para resíduos
mistos soltos e de 237 a 415~kg/m³ para resíduos compactados
\cite{epa2016volume}. A diferença entre essas fontes já mostra que converter
volume em massa exige hipóteses explícitas; por isso, o modelo deste trabalho
opera em volume e só converte para massa quando a fonte do fator está indicada.

\subsection{Cidades inteligentes e IoT na gestão de resíduos}

Na gestão de resíduos, as aplicações de IoT mais frequentes são o
monitoramento do nível de enchimento, a identificação de contêineres e veículos
e a roteirização da coleta \cite{anagnostopoulos2017challenges,esmaeilian2018future}.
Uma revisão sistemática com 173 estudos primários mapeou sensores, atuadores,
partes interessadas e dados trocados nesses sistemas \cite{sosunova2022iot}.
Outra revisão define a lixeira inteligente como o contêiner dotado de
monitoramento de nível e conectividade IoT e aponta a falta de padronização
entre as soluções \cite{zoumpoulis2024smart}. Quanto à comunicação, um
levantamento sobre sistemas baseados em sensores identificou o LoRaWAN como o
protocolo preferido para automação em áreas urbanas \cite{vishnu2022sensor},
o que contrasta com o Wi-Fi usado no protótipo deste trabalho. Revisões sobre
inteligência artificial no setor ressaltam que qualidade de dados, privacidade e
custo ainda limitam a adoção \cite{olawade2024smart}.

Uma arquitetura IoT para esse fim costuma ser descrita em camadas: dispositivos
de borda com sensores e microcontrolador, uma camada de comunicação e uma
camada de aplicação que armazena, processa e apresenta os dados
\cite{oliveira2018arquitetura}. O protótipo de 2024 seguiu essa organização.
A presente versão a desdobra em sete camadas, separando a ingestão, a
persistência, a inteligência operacional, a visualização e a avaliação
(Seção~\ref{sec:arquitetura}).

\subsection{Medição de nível por ultrassom}

O HC-SR04 emite um pulso de 40~kHz e mede o tempo até o retorno do eco. Sua
folha de dados indica faixa de 2 a 400~cm, ângulo de medição de 15° e precisão
de até 3~mm, e recomenda alvo com área mínima de 0,5~m² e superfície ``o mais
lisa possível'' \cite{elecfreakssdhcsr04}. A superfície de resíduos em um
contêiner não atende a essas condições: é irregular, pode ter sacos inclinados
e varia de material. Como o sensor registra o primeiro eco recebido dentro do
cone, a distância medida tende a corresponder ao ponto mais alto do monte, e
não ao nível médio. Um estudo exploratório comparou o monitoramento por sensores ultrassônicos com
a observação visual dos motoristas e analisou, para cada abordagem, a troca
entre o número de coletas e o número de transbordamentos
\cite{brouwer2023comparison}.

A velocidade do som no ar varia com a temperatura. A aproximação linear
$c \approx 331{,}3 + 0{,}606\,T$~m/s (com $T$ em °C) mostra variação de cerca
de 5\% entre 10 e 40~°C. O DHT11 mede temperatura de 0 a 50~°C com exatidão
típica de $\pm1$~°C \cite{drobotics2010dht11}, o que é suficiente para essa
compensação. No protótipo original, o DHT11 foi apresentado como um sensor que
influenciaria o nível de enchimento; nesta versão, seu papel limita-se à
compensação da velocidade do som e ao registro de condições ambientais.

\subsection{O protocolo MQTT}

O MQTT é um protocolo de mensagens do tipo publicação/assinatura padronizado
pela OASIS, nas versões 3.1.1 \cite{banks2014mqtt} e 5.0 \cite{banks2019mqtt}.
Os clientes não se comunicam diretamente: publicam mensagens em
\emph{tópicos} hierárquicos de um \emph{broker}, que as entrega aos clientes
com assinaturas compatíveis, inclusive por curingas (\texttt{+} para um nível,
\texttt{\#} para vários). Quatro recursos do protocolo são usados neste
trabalho:

\begin{itemize}
  \item \textbf{Qualidade de serviço (QoS).} Com QoS 0, a mensagem é enviada no
  máximo uma vez, sem confirmação; com QoS 1, ao menos uma vez, com
  confirmação e possível duplicação; com QoS 2, exatamente uma vez, por meio
  de um protocolo de quatro passos \cite{banks2014mqtt}. A garantia vale entre
  cada cliente e o broker; não há garantia fim a fim se um cliente perde dados
  antes de publicá-los.
  \item \textbf{Mensagens retidas.} O broker guarda a última mensagem retida de
  cada tópico e a entrega imediatamente a novos assinantes, o que é útil para
  o estado de disponibilidade de um dispositivo.
  \item \textbf{\emph{Last Will and Testament} (LWT).} Na conexão, o cliente
  registra uma mensagem que o broker publica se a conexão for encerrada sem a
  mensagem DISCONNECT, por exemplo após queda de energia ou expiração do
  \emph{keepalive}.
  \item \textbf{Sessões persistentes.} Com \emph{clean session} desativado, o
  broker mantém as assinaturas e as mensagens QoS~1/2 pendentes de um cliente
  desconectado, entregando-as na reconexão.
\end{itemize}

Essas garantias dependem do cliente usado. A biblioteca PubSubClient, empregada
no firmware do ESP32, implementa o MQTT 3.1.1, publica apenas com QoS~0 e não
mantém fila de mensagens durante desconexões (ver o relatório de auditoria do
firmware no repositório). Por isso, embora o broker Eclipse Mosquitto
\cite{light2017mosquitto} aceite o MQTT~5.0, todo o sistema foi especificado em
MQTT~3.1.1, e nenhuma afirmação sobre recursos do MQTT~5.0 é feita.

O MQTT não define autenticação nem cifragem: elas dependem do broker e do
transporte (TLS). Brokers abertos à internet sem autenticação permitem leitura
e injeção de mensagens por terceiros, e o protocolo é alvo de ataques de negação de
serviço de baixa taxa, como o SlowITe, demonstrado contra serviços MQTT reais
com e sem cifragem \cite{vaccari2020slowite}. Estudos comparativos analisam o desempenho do MQTT
frente a protocolos como o CoAP, inclusive com mecanismos de segurança ativados
\cite{seoane2021performance}.

\subsection{Modelagem, simulação e gêmeo digital}

Simulação de eventos discretos e simulação em tempo discreto são abordagens
usuais para avaliar sistemas logísticos estocásticos quando a experimentação
real é cara ou inviável \cite{law2024simulation}. A credibilidade dos
resultados depende de verificação (o modelo computacional implementa
corretamente o modelo conceitual) e de validação (o modelo conceitual
representa o sistema real com exatidão suficiente para o propósito)
\cite{sargent2013verification}. Quando várias alternativas são comparadas, o
uso de \emph{números aleatórios comuns}, isto é, a submissão de todas as
alternativas às mesmas realizações das variáveis aleatórias, reduz a
variância das diferenças estimadas e permite comparações pareadas
\cite{glasserman1992some}.

O termo \emph{gêmeo digital} é usado com frequência em plataformas urbanas de
monitoramento \cite{barth2023data,cardenasleon2024urban}, nem sempre com o
mesmo significado. A revisão de \cite{kritzinger2018digital} distingue, segundo
o nível de integração de dados entre o objeto físico e sua contraparte digital,
três categorias: \emph{modelo digital}, \emph{sombra digital} e \emph{gêmeo
digital}\footnote{Apenas o resumo desse artigo foi consultado; o texto completo
não pôde ser obtido por acesso automatizado.}. Neste trabalho, adota-se o
seguinte critério operacional, coerente com essa gradação: há sombra digital
quando o estado do objeto físico atualiza automaticamente o modelo, e só se
fala em gêmeo digital quando, além disso, as decisões do modelo retornam
automaticamente ao objeto físico. A plataforma desenvolvida mantém, para cada
contêiner, um estado estimado atualizado automaticamente pela telemetria, mas
as decisões voltam ao ``mundo físico'' por meio de rotas executadas por
veículos, e esse mundo é, ele próprio, simulado. A solução corresponde,
portanto, a uma \emph{sombra digital} acoplada a um simulador. Por esse motivo,
o termo ``gêmeo digital'', considerado no título provisório deste trabalho, não
foi adotado.

\subsection{Roteirização de veículos}

O problema de roteirização de veículos (VRP) foi formulado por
\cite{dantzig1959truck} e consiste em definir rotas de custo mínimo que partem
de um depósito e atendem um conjunto de clientes. Na variante capacitada
(CVRP), a soma das demandas de cada rota não pode exceder a capacidade do
veículo; restrições de duração de jornada e de janelas de tempo geram outras
variantes \cite{toth2014vehicle}. Como o CVRP é NP-difícil, instâncias
práticas são resolvidas por heurísticas: construtivas, como o vizinho mais
próximo \cite{rosenkrantz1977analysis} e as economias de
\cite{clarke1964scheduling}; de melhoria local, como o 2-opt
\cite{croes1958method}; e metaheurísticas, como a busca local guiada (GLS)
\cite{voudouris1999guided}. A biblioteca de roteirização do Google OR-Tools
combina essas técnicas em um resolvedor de propósito geral
\cite{furnon2025ortools,perron2025ortools}. Nenhuma dessas heurísticas garante
a solução ótima, e os resultados deste trabalho não são apresentados como
ótimos.

A coleta de resíduos é um caso particular bastante estudado
\cite{belien2014municipal,han2015waste,hess2024waste}. Quando o nível dos
contêineres é conhecido, a decisão diária passa a incluir \emph{quais}
contêineres visitar, além da ordem de visita. Essa estrutura aproxima o
problema do roteamento de inventário, no qual se decide simultaneamente quando
reabastecer (aqui, esvaziar) cada ponto e como roteirizar \cite{mes2014inventory}.
Uma formulação conveniente é o CVRP com prêmios (\emph{prize-collecting}): cada
contêiner não visitado gera uma penalidade, e o resolvedor equilibra a
distância adicional com o benefício de visitar. Como os volumes reais só são
conhecidos na chegada, a execução pode exigir um desvio ao depósito quando o
caminhão enche, o que corresponde à política de recurso clássica do VRP com
demandas estocásticas \cite{dror1989vehicle}.

\subsection{Avaliação experimental}

Em experimentos de simulação com réplicas independentes, cada semente aleatória
produz uma observação de cada métrica. Com números aleatórios comuns, as
observações de duas políticas na mesma semente formam pares. O teste de postos
sinalizados de \cite{wilcoxon1945individual} compara pares sem supor
normalidade das diferenças; quando várias comparações são feitas, o
procedimento de \cite{holm1979simple} controla a taxa de erro do tipo~I da
família de testes. Intervalos de confiança podem ser obtidos por reamostragem
\emph{bootstrap} \cite{efron1979bootstrap,efron1993introduction}, e a correlação
bisserial de postos pareada, calculada pela fórmula da diferença simples
\cite{kerby2014simple}, fornece um tamanho de efeito entre $-1$ e $1$. A
ausência de significância estatística não demonstra equivalência; apenas indica
que os dados não permitem distinguir as alternativas com a confiança adotada.
Por fim, a reprodutibilidade de experimentos computacionais exige registrar
versões de código e dependências, sementes, parâmetros e dados brutos
\cite{sandve2013ten,peng2011reproducible}.
